\documentclass[12pt]{article}

\usepackage{amsmath}
\usepackage{amssymb}
\usepackage{booktabs}
\usepackage{geometry}
\usepackage{setspace}

\geometry{margin=1in}
\setstretch{1.2}

\title{CSE400 -- Lecture Scribe\\
Joint Random Variables and Joint Distributions}
\author{Course: Fundamentals of Probability in Computing}
\date{Instructor: Dhaval Patel, PhD \\ Lectures: L15--L18}

\begin{document}

\maketitle

\section{Outline of the Lecture}

The lecture covers the following topics:

\begin{itemize}
\item Why study joint random variables
\item Discrete case
\begin{itemize}
\item Joint PMF
\item Joint PMF table representation
\item Example: Rolling two dice
\end{itemize}
\item Continuous case
\begin{itemize}
\item Joint PDF
\item Geometric interpretation
\item Worked example
\end{itemize}
\item Joint cumulative distribution function (Joint CDF)
\begin{itemize}
\item Definition
\item Continuous case
\item Discrete case
\item Properties of the joint CDF
\item Comprehensive example
\end{itemize}
\end{itemize}

\section{Why Study Joint Random Variables?}

In many real systems, multiple random variables occur at the same time. Instead of analyzing them individually, we often need to understand their relationship and combined behavior.

Thus, we study \textbf{Joint Random Variables}, which allow analysis of two or more random variables simultaneously.

\section{Motivating Examples}

\subsection{Example 1: Smart Transportation}

Consider two random variables:

\[
X = \text{Vehicle Speed}
\]

\[
Y = \text{Traffic Density}
\]

Question:

Can speed be high when traffic density is also high?

This question studies the relationship between the two variables.

Thus:

\[
X \text{ and } Y \text{ are correlated random variables.}
\]

Applications in Intelligent Transportation Systems:

\begin{itemize}
\item Adaptive traffic signals
\item Congestion prediction
\item Autonomous vehicle routing
\end{itemize}

\subsection{Example 2: Wireless Network Performance}

Two important random variables:

\[
X = \text{Signal-to-Noise Ratio (SNR)}
\]

\[
Y = \text{Data Throughput}
\]

Key reasoning:

Throughput depends on SNR.

However, both vary due to:

\begin{itemize}
\item fading
\item interference
\item congestion
\end{itemize}

Question:

If SNR is high, will throughput always be high?

Answer: Not necessarily, because network congestion can influence throughput.

Applications:

\begin{itemize}
\item 5G scheduling
\item Adaptive modulation
\item Network planning
\end{itemize}

\subsection{Example 3: Weather Prediction}

Random variables:

\[
X = \text{Rainfall}
\]

\[
Y = \text{Temperature}
\]

Application: Weather forecasting.

\subsection{Example 4: Drone Delivery}

Random variables:

\[
X = \text{Wind Speed}
\]

\[
Y = \text{Delivery Time}
\]

Application: Logistics and UAV systems.

\subsection{Example 5: Rolling Two Dice}

Let:

\[
X = \text{value of first die}
\]

\[
Y = \text{value of second die}
\]

Example questions:

\[
X + Y = 7
\]

\[
X > Y
\]

Such questions require the \textbf{joint probability distribution}, specifically the joint PMF.

\section{Discrete Case -- Joint PMF}

When the random variables are discrete, their joint behavior is described using the \textbf{Joint Probability Mass Function (Joint PMF)}.

\[
p_{X,Y}(x,y) = P(X = x, Y = y)
\]

This value gives the probability that:

\begin{itemize}
\item random variable $X$ equals $x$
\item random variable $Y$ equals $y$
\end{itemize}

at the same time.

\section{Joint PMF Table Representation}

A convenient way to represent the joint PMF is through a table format.

\begin{center}
\begin{tabular}{c|ccc}
\toprule
$Y \backslash X$ & $x_1$ & $x_2$ & $x_3$ \\
\midrule
$y_1$ & $P(x_1,y_1)$ & $P(x_2,y_1)$ & $P(x_3,y_1)$ \\
$y_2$ & $P(x_1,y_2)$ & $P(x_2,y_2)$ & $P(x_3,y_2)$ \\
\bottomrule
\end{tabular}
\end{center}

Each cell contains the probability:

\[
P(X=x, Y=y)
\]

Thus the table lists probabilities for every possible pair of values.

\section{Example: Rolling Two Dice}

Consider rolling two dice.

Define:

\[
X = \text{result of first die}
\]

\[
Y = \text{result of second die}
\]

Each die can take values:

\[
1,2,3,4,5,6
\]

Thus the total number of possible outcomes is:

\[
6 \times 6 = 36
\]

Each outcome is a pair:

\[
(X,Y)
\]

Examples of possible pairs:

\[
(1,1), (1,2), (2,5), (6,6)
\]

For fair dice:

Each pair has probability

\[
P(X=x,Y=y)=\frac{1}{36}
\]

Using this joint PMF table, we can compute probabilities of events such as:

\[
X+Y=7
\]

\[
X>Y
\]

\section{Continuous Case -- Joint PDF}

When random variables are continuous, we use the \textbf{Joint Probability Density Function (Joint PDF)}.

\[
f_{X,Y}(x,y)
\]

This function represents the density of probability for the pair $(X,Y)$.

\section{Geometric Interpretation}

For continuous variables:

The joint PDF can be visualized as a surface over the $x$-$y$ plane.

The probability corresponds to the \textbf{volume under the surface} over a given region.

Thus the probability for a region is obtained by integrating the joint PDF over that region.

\[
P((X,Y)\in R)=\int\int_R f_{X,Y}(x,y)\,dx\,dy
\]

\section{Worked Example (Continuous Case)}

In solving continuous joint probability problems, the following steps are used:

\begin{enumerate}
\item Identify the region where the joint PDF is defined.
\item Express the probability event that needs to be calculated.
\item Set up the double integral over the region.
\item Integrate step by step.
\item Obtain the required probability.
\end{enumerate}

This demonstrates how probabilities are computed for continuous joint random variables.

\section{Joint Cumulative Distribution Function (Joint CDF)}

The Joint CDF describes the cumulative probability for two random variables.

\[
F_{X,Y}(x,y)=P(X\le x, Y\le y)
\]

This represents the probability that:

\begin{itemize}
\item $X$ is less than or equal to $x$
\item $Y$ is less than or equal to $y$
\end{itemize}

\section{Joint CDF -- Continuous Case}

For continuous random variables, the joint CDF is obtained by integrating the joint PDF.

\[
F_{X,Y}(x,y)=\int_{-\infty}^{x}\int_{-\infty}^{y} f_{X,Y}(u,v)\,du\,dv
\]

This accumulates the probability from negative infinity up to $x$ and $y$.

\section{Joint CDF -- Discrete Case}

For discrete random variables, the joint CDF is obtained by summing the joint PMF values.

\[
F_{X,Y}(x,y)=\sum_{u\le x}\sum_{v\le y} P(X=u,Y=v)
\]

This means we add the probabilities of all pairs satisfying:

\[
X \le x \quad \text{and} \quad Y \le y
\]

\section{Properties of the Joint CDF}

\subsection{1. Range}

\[
0 \le F_{X,Y}(x,y) \le 1
\]

since it represents probability.

\subsection{2. Non-decreasing behavior}

The joint CDF does not decrease as $x$ or $y$ increases.

\subsection{3. Limit behavior}

As $x$ and $y$ approach infinity:

\[
F_{X,Y}(x,y)\rightarrow 1
\]

because all probability mass is accumulated.

\section{Key Points for Exam Revision}

\subsection{Why Joint Random Variables?}

Used to analyze two or more random variables together.

\subsection{Discrete Case}

Joint PMF:

\[
p_{X,Y}(x,y)=P(X=x,Y=y)
\]

Represented using joint PMF tables.

\subsection{Continuous Case}

Joint PDF:

\[
f_{X,Y}(x,y)
\]

Probabilities computed using double integrals.

\subsection{Joint CDF}

\[
F_{X,Y}(x,y)=P(X\le x,Y\le y)
\]

\begin{itemize}
\item Continuous case $\rightarrow$ integration
\item Discrete case $\rightarrow$ summation
\end{itemize}

\end{document}